\chapter{Focus and focus realization} \label{focusandfocusrealization} \label{sec:1}

Following \textcite{zimmermann2011focus}, I consider focus a universal category within the framework of Information Structure (hereafter IS). I also distinguish between the semantico-pragmatic category of focus and its grammatical manifestation in a given language.
In other words, I argue that the notion of focus should be separated from its formal grammatical realization in a sentence, i.e., the syntactic domain where it is expressed and the prosodic or morphological means by which it is marked. Accordingly, I disagree with theories that define focus through specific accents or dedicated syntactic positions (e.g., \cite{chomsky1970some, rizzi1997fine, brody1990some}). Instead, I assume that even in the absence of special marking, the informative value of an utterance is still conveyed.

In the first part of Chapter 1, the concept of focus is discussed independently of how it is marked, with an emphasis on its cross-linguistic function and its contribution to the universe of discourse. In line with \textcite{roberts1996information} and others, I view focus as a discourse-structuring device. Specifically, the central function of focus in declarative utterances is to identify the questions these utterances can answer \parencite{roberts1996information, buring2003d, beaver2009sense}.
The realization of focus, addressed in the second part of the chapter, refers to its linguistic expression through special grammatical means, such as prosodic prominence, syntactic reordering, and morphological markers. These devices help the hearer identify the focus (and, if possible, the background) structure of an utterance.
This chapter is organized as follows: In Section 1.1, I provide a non-exhaustive synthesis of various contributions to the notion of focus as an information-structural category. I then clarify the distinction between focus structure and different focus types (Sections 1.1.1 and 1.1.2). Section 1.2 presents the different ways languages mark focus. Finally, Section 1.3 concentrates on subject focus, offering a brief overview of its syntactic manifestations across languages.

\section{The notion of Focus} 

The notion of focus has its origins in discourse structure, semantics, and pragmatics  (\cite{krifka2008basic, rooth1985association, rooth1992theory, schwarzschild1999givenness, von1990current}, \textit{inter alia}). Given the complexity of the phenomenon, it is not surprising that the concept has been subject to a long and intricate debate over the years. 
It is thus difficult to trace back the history of focus as a linguistic notion, because the concept was expressed in different terms, and often the same term was used to express different notions. The Prague school (e.g., \cite{mathesius1983functional, danevs1964three, firbas1971concept}), for instance, made a distinction between \textit{theme}, that part that refers to what is already known or given in the context (also known by other scholars as \textit{the topic}, see \cite{reinhart1981pragmatics, jacobs2001dimensions}, among others) and \textit{rheme}, the part that conveys new information (which corresponds to the focus of the sentence). In the sentence \textit{John bought a new car}, for instance, \textit{John}, the referent already known or given in discourse, was considered the \textit{theme}, while the sentence \textit{bought a new car}, which predicates something new about John, was called \textit{rheme}.

In generative grammar, the term focus was used by \textcite{chomsky1970some, jackendoff1972semantic,brody1990some, rizzi1997fine}, who consider focus to be the information that the speaker assumes is not shared by the hearer, while non-focused material conveys the \say{presuppositional skeleton}, a piece of information already shared by both speaker and hearer \parencite[2122]{frascarelli2010narrow}.  \textcite{jackendoff1972semantic} introduced a syntactic feature named F. The F feature marks the focused phrase, while a phrase which is not marked with F is unfocused. F would correspond to stress prominence in a certain domain. 
 
In a similar vein, the Cartographic approach \parencite{rizzi1997fine}, building on the observation that focus is marked in several languages by various grammatical means, assumes focus to be a part of the syntactic structure, at least in the form of an abstract feature, located in a dedicated position in the C-domain (the left periphery of the sentence or LP). This is known as the Focus Phrase or FocP:

\ea
 
{} [ForceP [TopP* \textbf{[FocP XP[+foc]} [TopP* [FinP [IP ...tXP]]]]]]\ \\
\z


Thus, scholars from the generative tradition claim that IS-notions like topic and focus represent distinct structural positions in the left periphery of the sentence above TP, and that such notions are defined in terms of abstract features, which trigger specific marking-strategies.

Unlike the Cartographic approach, \textcite{matic2013meanings} argue that focus is not a core grammatical category, but rather a by-product of other linguistic systems such as aspect and modality. Indeed, even though focus-marking strategies are widely found cross-linguistically, there exist languages that do not mark focus at all.

\textcite{buring2009towards}, in his typology of focus realization, also discuss a series of languages that do not show any particular focus-marking strategy (e.g Hausa, see \cite{green2003ex, hartmann2008focus}). 

The lack of focus realization does not imply, however, that the concept is not present in the language.
According to \textcite{zimmermann2011focus}, formal marking of focus is not expected to be mandatory in each and every language. Languages do not necessarily make use of focus-marking devices that would yield a strict one-to-one mapping between focus and its grammatical realization. In their words, ``focus is expected to be formally underspecified, and thus in need of contextual resolution, in many, if not all languages. In fact, since the utterance context heavily reduces the number of worlds relative to which an utterance is evaluated [...], the explicit realization of focus might even seem superfluous at times” \parencite[1658]{zimmermann2011focus}.

Adopting this view, I consider that defining focus only by its grammatical, prosodic or morphologic realization is misleading, and I assume that focus is better identified via pragmatics. Concretely, I consider that certain (semantic or pragmatic properties of certain) contexts systematically trigger focus.

In a dialogical context, the classical pragmatic use of focus is to highlight the part of an answer that corresponds to the wh-part of a constituent question \parencite{paul2010prinzipien}. \textcite{drubig2003toward} claims that a wh-question always requires a focused answer. Consequently, answers to wh-questions are commonly used as diagnostics for focus and its grammatical realization. Consider the example in (\ref{ex12}):

\ea \label{ex12}
\ea Who ate the cake? 
\ex{} [Mary]\textsc{foc} \normalsize ate it.
\z \z

In earlier investigations, focus has often been associated with concepts such as \say{information point} of the sentence \parencite{bolinger1985two}, \say{emphasis}, \say{importance}, \say{highlighting} , and \say{newness} \parencite{halliday_1967}. However, these concepts have been considered unsatisfactory and have failed to capture the complexity of the phenomenon in its entirety. The notion of newness, for instance, clearly gives wrong predictions. Consider the case in which a constituent that refers to something previously mentioned is in focus, such as in the case of focused subjects. Subjects are usually entities that have already been previously introduced and, as such, are discourse-given. However, a subject can at the same time represent the focus of the sentence.

An important clarification in this respect is provided by \textcite{lambrecht1996information}, who considers focus (like topic) as a relational pragmatic category, that has to be distinguished from a referential category. While a relational category expresses a relation between the referent and the corresponding proposition, a referential category reflects the identifiability and activation states of discourse referents. Pronouns, for instance, which are the most activated referents, can also represent new information to the hearer, because they are in a relation of newness with respect to the proposition they refer to. Thus, in example (\ref{ex12}), Mary, being a proper name, has a uniqueness reference, and is already known by both speaker and hearer. Yet, the information that it was her who ate the cake is new to the hearer. 

It has been customary to draw a sharp division between semantic and pragmatic uses of focus. The classic example of semantic uses of focus is its association with exclusive particles like `only’. In general, ‘A only B’ means that, out of some range of properties A could have, the single property A actually has is B. Other alternatives are thus excluded (hence the term ‘exclusive’) \parencite{velleman2016question}. For instance, a sentence like `Peter only dances' indicates that the only property that Peter has is to dance (and not, say, to sing).
The semantic use of focus affects the truth-conditional content of the CG. In this case, failing to set the focus right will result in the transmission of unintended factual information. Thus, the variation in truth conditions between these sentences shows that focus helps determine which alternatives are excluded.

\textcite{rooth1985association}, in his Alternative Semantic Framework \parencite{rooth1985association, rooth1992theory}, describes focus as indicating alternatives, and a focus sensitive particle is one whose meaning depends on those focal alternatives.
Such truth-conditional effects of focus can be distinguished from the `pragmatic’ uses of focus. These uses are connected to the public communicative goals of the speakers, and do not have an immediate influence on truth conditions. Rather, they contribute to guide the direction in which communication should develop. This implies that failing to select the right focus typically results in incoherent communication. 

In subsequent research, it was shown that this unconstrained system of Alternative Semantics makes available too many alternatives, and allows for a number of incorrect predictions. To solve this issue, \textcite{rooth1992theory} proposes a more restrained version of Alternative Semantics, by adding a context variable C, which refers to a contextually-given set of alternatives, and a focus operator, which introduces a presupposition requiring the value of C to be a subset of the (unconstrained) set of focus alternatives.

While in Alternative Semantics focus only helps to identify the relevant alternatives, which are registered at an independent level of meaning, the Structured Meaning approach \parencite{von1982structured, von1990current, jacobs2012fokus, krifka1991compositional, krifka2008basic} assumes that focus leads to a partition of the sentence meaning into two parts: focus and background. In the Structured Meaning approach, focus has the effect of structuring the propositions denoted by sentences: the focus-influenced semantic value of a clause with a single focus is a pair consisting of (i) a property, obtained by lambda-abstracting on the focused position, and (ii) the semantics of the focused phrase.
In this respect, \textcite[38]{von1991focusing} states that the semantics of questions is fundamental for a focus theory: a wh-question sets the background for an answer, which, in turn, determines the focus of the answer. 

A theory that takes questions as the ``touchstone” for focus is the Questions Under Discussion (QUDs) approach (see, for instance, \cite{roberts1996information, buring2003d, beaver2009sense}). Within this theory, 
questions, explicit or implicit, are seen as structuring discourse, and information-structural marking is seen as reflecting that underlying discourse structure. Moreover, this account collapses the distinction between semantic and pragmatic uses of focus, and analyses all effects of focus (including semantic ones) as arising from a single essentially pragmatic function: focus helps to indicate which QUD is the current question, the question the current discourse move is intended to address.

QUDs are taken to represent discourse goals, and focus-marking in an assertion indicates which QUD it answers (which question is the current question with respect to the assertion), thereby indicating the discourse goal that needs to be fulfilled.

QUDs are often considered as sets of alternative propositions and are therefore related to \textcite{rooth1985association, rooth1992theory}'s Alternative Semantics. Indeed, according to \textcite{velleman2016question}, the question-based approach to focus can be treated as an add-on to Alternative Semantics rather than a replacement.
However, while Alternative Semantics says little about how listeners choose the right logical form, except in cases where the discourse makes an antecedent set of alternatives explicit, the QUD approach sets up a series of constraints which are claimed to enable speakers to arrive at the right focal alternatives. In other words, when speakers hear a sentence whose information structure is ambiguous, they would choose the structure that leads to the fewest constraint violations and the least need for accommodation or repair.
Thus, this approach is not incompatible with the Alternative Semantic Framework and Structured Meaning, but rather builds on previous claims about the theory of focus. Like the other two approaches, within the QUD framework, focus is defined as the element in the answer which instantiates the open variable in the QUD. 

The idea that the focus constituent is the part of the sentence that corresponds to the answer to a question \parencite{bresnan1989locative} is not new. Within this framework, however, an answer to a question is appropriate only if its focused constituent corresponds to the wh-phrase of the question. 
Importantly, the QUD approach posits a similar hierarchical structure even for discourses which do not involve any explicitly asked questions, such as in corrections (contrastive or corrective focus) or in out-of-the-blue utterances, in which case they can only be identified with recourse to the focus alternatives.
 In these cases, the QUD has to be reconstructed starting from the answer itself and looking at the preceding discourse context. To reconstruct an implicit QUD, \textcite{riester2018annotation, riester2019constructing} propose annotation guidelines for QUD, which will be discussed in detail in the methodology section in Chapter 4. 
 
The QUD approach to focus presents several methodological advantages. First of all, as mentioned above, it is compatible with both the Alternative Semantics and the Structured Meaning approach. Secondly, it provides specific indications and constraints on how to retrieve the correct information-structural partition of the sentence and the consequent identification of focus. Third, it clearly indicates how to reconstruct the QUD when this is implicit, thus avoiding exclusive reliance on the syntactic form of the answer, which, as will be shown, is often underspecified and can be compatible with more than one sentence partition. Given the empirical nature of the book, in that it builds on dialogical data, implicit questions are not uncommon, and the QUD method proves particularly useful in this regard.
 
Finally, the correct retrieval of the QUD allows one to distinguish between focus types with a contrastive and non-contrastive interpretation. In the following sections, I will discuss the contributions to different focus structures and different focus types, and provide some examples of how the QUD method can help to differentiate efficiently among different sub-categories.

\subsection{Focus structure}
As is generally acknowledged, constituents of different \say{sizes} can be focused within the sentence. One of the major contributions to the focus structure of the sentence is provided by \textcite{lambrecht1996information}, who defines focus structure as the conventional association of a focus meaning (distribution of information) with a sentence form. 

A fundamental contrast is observed between broad and narrow focus \parencite{selkirk1984phonology, lambrecht1996information}. Such a distinction is not categorical but gradual in nature, with narrow focus referring to focus on single words or constituents and broad focus referring to complex constituents (VP) or to the case in which all parts are newly introduced into the discourse at the moment of utterance \parencite{fery2008information}.
\textcite[611]{lambrecht2000subjects} distinguishes among Argument Focus (AF), Predicate Focus (PF) and Sentence Focus (SF), which are described as follows:

\begin{description}
  \item[\textbf{Sentence focus structure} (SF),] or all-focus sentence, expresses a pragmatically structured proposition in which both the subject and the predicate are in focus. The focus domain is thus the whole sentence. This sentence type differs from predicate focus constructions, in that it exhibits no topical element, and is felicitous as an answer to questions like ‘What happened?'. Consider the example in (\ref{ex13}):
    
\ea \label{ex13}
\ea What happened?
\ex \ [My car broke down]\textsc{foc}\normalsize. \hfill \parencite[223]{lambrecht1996information}   
\z \z

\item[\textbf{Predicate focus structure} (PF),] also called VP focus, is considered universally the unmarked type by \textcite[622]{lambrecht2000subjects}, and coincides with the traditionally recognized topic-comment organization of information in a sentence. The Predicate focus structure is a sentence construction expressing a pragmatically-structured proposition in which usually the subject (although it can also be another constituent) is the topic (hence within the presupposition) and in which the predicate expresses new information about this topic. The focus domain is, thus, the predicate phrase (or part of it). An example is provided in (\ref{ex14}):
    
    \ea \label{ex14}
    \ea What did Laura do yesterday?
    \ex She [went to the park]\textsc{foc}\normalsize.
    \z \z

\item[\textbf{Argument Focus structure} (AF)] refers to the situation in which the focus domain is represented by a single constituent that has an argumental status (\ref{ex15}):

\ea \label{ex15}
\ea Who stole the book?
\ex \ [Marc]\textsc{foc} \normalsize stole it.
\z \z

\end{description}


The three focus structures are used in different communicative situations, and are felicitous as answers to different QUDs. As mentioned above, \textcite{lambrecht2010constraints} considers the Predicate focus (PF) to be the \say{unmarked} structure and the one most frequently found across languages. 
Conversely, the Argument focus (AF) category, where the focused constituent is the subject (\say{identificational} sentence in Lambrecht's terms), is described as the most \say{marked}, infrequent case. 
In this book, I will concentrate mainly on the Argument Focus structure, this being the category to which narrowly-focused subjects belong. The other two categories and their mapping onto grammar are beyond the scope of the present work.
In the Argument focus condition, the focused constituent can trigger different pragmatic inferences. These interpretations correspond to different sub-types of focus, and are discussed in the next section.

\subsection{Focus types} 

A considerable amount of work, from both a generative and functional perspective,  has been dedicated to distinguishing different types of foci, relying mainly on semantic and pragmatic differences. Due to the complexity of the phenomenon, the terminology used by different linguists has often been inconsistent, which has prevented, at times, a reliable comparison of the interpretation of certain kinds of focus across languages. 
There is a general tendency, advocated mainly by \textcite{bolinger1972intonation} and \textcite{schmerling1973aspects}, to distinguish between ``neutral” focus, whose main function is to express the most important part of the utterance (see, for instance (\ref{ex15})), or what is new in the utterance, and other pragmatic uses of focus that aim to correct information (also called `contrastive', `corrective' or non-neutral).\\
According to \textcite{krifka2008basic}, in the case of ``contrastive/corrective” focus, the focus alternatives must include a proposition that has been proposed in the immediately-preceding CG. This leads to a corrective interpretation in case the context proposition was different (see also \cite{truckenbrodt1995phonological, rooth1996focus, steube2001correction, van2004incompatibility, REPP20101333} for an exhaustive analysis of this focus type).

 \textcite{kiss1998identificational} advocates a further distinction between identificational focus and information focus (see also \cite{rochemont1986focus}; \cite{campos1990focalization}). The two sub-types are  differentiated for semantic reasons: identificational focus performs exhaustive identification of a set of entities given in the context (discourse-given) or in the situation (situationally-given), while information focus simply marks the non-presupposed nature of the information it carries \parencite{riester2013focus}. Consider the following example:
 
\ea \label{ex19}   {Identificational focus}\\
 \ [\textit{PÉTER}]\textsc{foc} \normalsize \textit{aludt a padlón}.\\
‘It was Peter who slept on the floor.’ \hfill \parencite[75]{kiss2008topic}   

\z

 \textcite{szabolcsi1981semantics} describes the exhaustive interpretation of the focus in (\ref{ex19}) with the following formula ``For every x, x slept on the floor if x = Péter”. Thus, identificational focus excludes, by definition, all the other possible alternatives of the set. In the case of information focus, on the other hand, such interpretation is not conveyed:

\ea \label{ex20}  {Information Focus}
\ea What did you bring Mary?
\ex I brought her [flowers]\tiny FOC.
\z \z

While in (\ref{ex19}) the only possible referent for which the proposition holds is Peter, such an implication is not conveyed in (\ref{ex20}). In other words, adding ‘and John too' to (\ref{ex19}) would be pragmatically infelicitous, while adding ‘and chocolates too' to the (\ref{ex20}) would be quite acceptable in context. As mentioned above, this would be attributable to the fact that identificational focus always carries an exhaustive interpretation, while information focus simply communicates the non-presupposed part of the information.

\textcite{kiss1998identificational} provides arguments in favour of such differentiation based on examples from Hungarian, where this distinction is said to be reflected in the grammar of the language. According to the author, identificational focus, being the most marked type, always triggers reordering of the constituents in the sentence (e.g., the use of a cleft sentence), while for information focus, which acts as a default type, such an operation is not required.

While Kiss gives a fundamental semantic distinction between focus types, the concept of exhaustivity is still a matter of debate. In particular, for some authors (see for instance \cite{horn1981exhaustiveness, destruel2018interpretation}), reordering of the constituents is not a reliable criterion to detect the exhaustive interpretation of the focus constituent. In some cases, an exhaustive focus can also be expressed by \textit{in situ} strategies \parencite{gabriel2010focus, fery2001focus, uth2014spanish}, and non-exhaustive focus can also involve reordering of the constituents. Rather, the context of the sentence, when possible, should be the only reliable way to determine exhaustiveness. 

\textcite{destruel2012french}, for instance, showed that associating additive particles such as `too' with \textit{it}-cleft sentences in French does not lead to ungrammaticality nor to pragmatic infelicity. These authors suggest that exhaustiveness is a pragmatic implicature (hence cancellable) rather than a semantic implication. The notion of exhaustivity associated with French clefts will be further discussed in Chapter 2 and Chapter 5 of this book.

Moreover, \textcite{kiss1998identificational} considers contrastive/corrective focus to be a sub-type of the identificational one.
Also according to \textcite{tomioka2008information}, the notion of contrastivity is connected to diverse linguistic phenomena, such as, e.g., exhaustive answers in question-answer pairs, contrastive statements, or instances of corrective focus. 

Thus, the term `contrastive' has also been commonly used to refer to cases where the choice of an element among a set of explicitly mentioned alternatives automatically excludes the other possibilities. Such exclusion can be carried out via two processes: either selecting one alternative among a set of explicit possibilities \parencite{dik1981typology}, or correcting a previously mentioned value. This is the position taken by \textcite{chafe1976givenness}, where contrast is thought to be any instance of focus where one or more alternatives to the focused expression have been explicitly mentioned in the preceding discourse. Moreover, it always carries an exhaustive interpretation. 

In some theories, the notion of contrast and correction are collapsed and reduced to cases where focus implies removal of information and insertion of new information. The two examples below exemplify the two semantic processes:

\ea \label{ex21}  {Contrastive Focus}
\ea Who stole the book, Marc or John?
\ex \ [Marc]\textsc{foc} \normalsize stole it.
\z \z 

\ea \label{ex22}  {Corrective focus}
\ea John stole the book
\ex  \ [Marc]\tiny FOC \normalsize{stole it (and not John)}.
\z \z

While some scholars advocate a further distinction between contrastive and corrective, \textcite{krifka2008basic} states that the distinction between identificational and information focus is unmotivated, and that the two categories should simply be combined under the definition of information focus (see, for instance, \cite{vallduvi1992informational}).
\textcite{krifka2008basic} claims that all focus types (contrastive or identificational or information focus) are in some sense tied to the notion of alternatives, and that the differences between the various interpretations can be captured in terms of the ways in which the ordinary meaning of the focused constituent interacts with the set of alternatives. Therefore, in his theory, he aims to simplify the differentiation among sub-types into  two major categories, namely \say{open} focus (open set of alternatives) vs. \say{closed} focus (restricted set of alternatives), and to restrict the term ``contrastive” only to cases which truly involve correction or additive particles (e.g., too).

Conversely, \textcite{torregrossa2012towards}, \textcite{cruschina2012discourse, cruschina2021,cruschina_mayol_2024_syntax_mapping, cruschina_mayol_2025_gradience_optionality} are increasingly calling into question the usefulness of the ``traditional” dualistic distinction between \textit{neutral} (information/identificational) focus and \textit{non-neutral} (contrastive, corrective, etc.) focus by providing evidence in favor of more fine-grained and gradual pragmatic distinctions.

Nevertheless, the distinction between \textit{neutral} and \textit{non-neutral} focus type continues to be particularly advocated due to its presence in the syntax of certain languages. For instance, it is commonly assumed that contrastive (hence non-neutral) foci licence certain word orders that would be judged infelicitous if they were to simply express information focus. A classic example is fronted foci in Spanish or Italian, where the common locus for focus is the postverbal slot of the sentence. Indeed, it is claimed that, if a focused constituent appears preverbally, it almost always has to bear a contrastive (or corrective) interpretation and a specific prosodic contour \parencite{beninca1988ordine, bianchi2013focus, zubizarreta1999word}. It must be noted, however, that there is no consensus in the literature on the question of specific focus types, and syntactic position or syntactic construction. In principle, contrast can also be encoded postverbally, while information can also be expressed \textit{in situ} \parencite{cruschina2021, cruschina_mayol_2025_gradience_optionality}. \textcite{brunetti2004}, for instance, argues that a fronted focus can simply answer a question. The question, however, has either to be implicit or not present in the immediately preceding discourse.

On a more general level, many linguists interested in the phenomenon  have investigated whether the semantic or pragmatic inferences observable across focus types and focus structures are in fact reflected in the grammar of natural languages.
By today, there seems to be a general tendency to answer this question in the negative. Although some patterns are visible, the strict one-to-one mapping between focus type and corresponding formal marking across languages remains largely theoretical. Instead, many structures appear underspecified with respect to the type of focus they encode, and are often compatible with different focus interpretations. This is  true for both focus structure and focus realization. The issue of realization will be discussed in the next section.

\section{Focus realization}
 The term focus realization, or focus marking, is generally used to indicate the formal mechanism that signals a focus relation between a pragmatically construed referent and a proposition \parencite{lambrecht1996information}, or to refer to the linguistic realization of focus by way of special grammatical means, such as, focus accenting, syntactic reordering, morphological markers, etc. \parencite{buring2009towards}.
Given the distinction between focus and focus realization, and given the above characterization of focus as a universal cognitive or semantic category at the level of information structure \parencite{zimmermann2011focus}, it is not the semantico-pragmatic category of focus, but rather its grammatical realization which is subject to cross-linguistic variation. 

As has previously been observed, languages show a wide range of grammatical means for marking a focused constituent as such, ranging from phonological and prosodic devices (accenting, prosodic phrasing), through morphological devices (focus markers), to syntactic devices (constituent reordering).
However, as mentioned above, focus marking is not mandatory in all languages. Rather, it seems to be formally underspecified (and therefore in need of contextual resolution) in many, if not all languages \parencite{zimmermann2011grammatical}. 

Thus, the grammatical realization of focus is often ambiguous Thus, the grammatical realization of focus is often ambiguous, in that it is compatible with more than one focus–background structure. Moreover, the term \textit{focus marking}, when used in this general way, implies that the formal devices used in the marking of focus are found exclusively in focus constructions \parencite{fiedler2006subject}.
Nevertheless, some cross-linguistic patterns can be identified. The range of possibilities for marking focus is rather wide, and it would be difficult to mention all the possible strategies. \textcite{lambrecht1996information} identifies the following classification: 

\begin{itemize}
    \item Exclusively prosodic (e.g., English, German)
    \vspace{-0.3cm}
    \item Morphological (e.g., Wolof, West Chadic)
     \vspace{-0.3cm}
    \item Prosodic and morphological (e.g., Japanese)
     \vspace{-0.3cm}
    \item Prosodic and syntactic (e.g., Italian, Spanish, French)
     \vspace{-0.3cm}
\end{itemize}

 In the generative account, the different focus-marking mechanisms have been regrouped into two main syntactic strategies, the \textit{in situ} and the \textit{ex situ} strategies \parencite{aboh2008focus}. The \textit{in situ} option is generally considered to be the most economical strategy, since the focused constituent is realized in the position defined by the argument structure of the verb, while languages that show the \textit{ex situ} option realize the focused constituent in the left periphery of the sentence (or C-domain in \cite{rizzi1997fine}), immediately followed by the verb. 
\textcite{buring2009towards}, in his Prominence Theory of Focus realization
\parencite{truckenbrodt1995phonological, buring2009towards}, claims that all focus-marking strategies, whether syntactic or not, can be analyzed under the umbrella of structural prominence, and that all languages are subject to one single constraint in focus realization, namely (\ref{ex23}):

\ea \textbf{FocusProminence}\\ \label{ex23}

Focus needs to be maximally prominent.
\z


 In Büring's account, the notion of \say{prominence} is proposed as a prosodic one (as originally intended by \cite{truckenbrodt1995phonological}). FocusProminence is thus considered to be the only way focus interacts with the formation of structure. Even in morphological focus-marking systems, morphemes can serve to mark heads and/or boundaries. In other words, these morphemes would not literally mark the focus but rather the structural position in which focus is realized.
However, adjustment strategies need not be prosodic: \say{if the prosodic system of a language does not allow A in the string A-B to be prominent, two more options arise: no adjustment whatsoever or syntactic adjustment to get A into a prosodically prominent position. The latter case will result in a language that marks focus by constituent order variation} \parencite[7]{buring2009towards}, as is the case in Spanish \parencite{zubizarreta1998prosody}. Given that FocusProminence makes reference to prosodic structure, this implies that, according to this view, marked constituent order is established only after prosodic structure is built.

Most of the semantic literature on focus assumes that a focused constituent contains a pitch accent, though focus may extend beyond the accented constituent itself (e.g., \cite{jackendoff1972semantic, selkirk2008contrastive}, \textit{inter alia}). Although such accents do not always affect the truth-conditional meaning of a sentence, especially in the absence of focus-sensitive operators like `only', they can strongly affect the appropriateness of an utterance \parencite{horn1981exhaustiveness, beaver2009sense}. \textcite{lambrecht1996information} agrees that the common denominator among different focusing strategies is the prosodic prominence of a given syllable in the sentence. It is also the only device which can occur by itself, without being complemented by another coding system (as is the case in English). It would seem therefore that the role of prosody in focus marking is, in some sense, functionally more important than morpho-syntactic marking. This is no doubt a consequence of the iconic relationship between pitch prominence and the degree of communicative importance assigned to the focal portion of a proposition. 

In French for instance, differences in the prosodic structure of a language account at least in part for the use of particular focus-marking systems. For example, it seems likely that the use of cleft constructions for the marking of focus differences is generally attributed to the fact that this language has both a relatively rigid constituent order and a relatively rigid rhythmic structure \parencite{dufter2003rhythmic}.\footnote{Note that the authors do not imply that cleft sentences are the only means to mark different focus types in French, and this claim is also supported by the corpus study reported in Chapter 3 of this book, where I show how different focus types can be expressed by several means in spoken French, albeit with different frequencies.}

Phenomena of non-canonical word order are particularly visible across languages in the case of contrast or correction. As argued in \textcite{zimmermann2008contrastive, REPP20101333}, the notion of contrast refers to the fact that a particular focus content, or a particular speech act containing a focus, or a particular focus-background partition, is unexpected for the hearer from the speaker’s perspective, and may thus create problems for the successful update of the CG. Since unexpected facts are more difficult to accommodate, the speaker will often try to facilitate the hearer’s task of adjusting her background assumptions accordingly for successful communication. One possibility for the speaker to direct the hearer’s attention, and to facilitate the task of shifting the background assumptions, is to use a non-canonical, i.e., a structurally more complex sentence that comes with additional grammatical marking in form of, for instance, a particular intonation contour, syntactic movement, a cleft structure, or the insertion of morphological markers. As \textcite{krifka2012expression} point out, there is a relatively strong need to express contrast cross-linguistically, and to express it with more explicit means such as syntactic movement. This can be explained by the plausible principle that if something is unexpected, it should be marked explicitly. 

If contrast and correction seem to be consistently marked, information focus-marking, on the other hand, often appears undetermined. This observation is particularly relevant in the case of object focus. As already mentioned in the section 1.1, object focus is consistently expressed across languages through \say{unmarked} strategies (see the ``Argument Asymmetry” in focus realization, \cite{skopeteas2010focus}). In languages that mark focus by means of a combination of prosodic and syntactic strategies, such as French and Spanish, realizing focus on the object generally results in the canonical SVO order. In the absence of explicit focus marking, the focus-background structure is massively underspecified by the grammatical system, and is hence in need of contextual resolution. Conversely, as has  often been stressed, a focus on the subject seems to require special marking across languages. As happens in the case of contrast or correction, it is often assumed that subject focus is explicitly marked because it conveys, once again, an \say{unexpected} interpretation of the subject as new information, that is more difficult to accommodate \parencite{lambrecht2000subjects}.
This subject/object asymmetry in focus realization will be discussed in detail in the next section.


\section{The Argument asymmetry in focus realization}

Previous work on IS has identified a well-known asymmetry with respect to the realization of focused constituents across numerous languages. 
It has already been observed for some languages that focus on subjects obligatorily induces a non-canonical structure, while focus on non-subjects only optionally does so \parencite{skopeteas2010focus}. This is presumably attributable to the frequent correlation between certain information-structural functions and certain syntactic constituents, where subjects are generally considered topics \say{by default}, while objects usually convey new information to the hearer.

Our utterances tend to be construed around and about subjects, which exhibit a range of semantic and discourse-related properties (e.g., givenness, agentivity, activation, etc.) which overlap with the ones displayed by topics. A subject in focus is, therefore, a marked option of the language and, consequently, less frequent. Given the general tendency in language to express marked interpretations through marked forms, it is not surprising that subject focus is consistently encoded by non-canonical strategies, while object focus generally lacks formal marking.

 Evidence for this subject/non-subject asymmetry in focus realization has been provided for several languages including Hausa \parencite{hartmann2007place}, West Chadic languages \parencite{zimmermann2011focus}, several Kwa and Gur languages \parencite{fiedler2005out} or Northern Sotho \parencite{zerbian2007investigating}.
The guiding hypothesis is that the special status of focused subjects is conditioned by information-structural factors, i.e., subjects in their canonical sentence-initial position are prototypically interpreted as topics \parencite{thompson1976subject}. Consequently, if the contextual cues establish a subject as focus, this conflicts with its primary information-structural function as non-focal information. In order to resolve this conflict, the focused subject will have to be realized in a non-canonical structure, for instance, by means of special morphological markers and/or syntactic reorganization \parencite{fiedler2006subject}. 

According to \textcite{lambrecht2010constraints}, speakers have an unconscious inclination to interpret isolated sentences as information-providing topic-comment sentences. Consequently, he considers that such articulation of the sentence represents, communicatively speaking, the most common type. In other words, it is more common for speakers to convey information about given discourse entities than to identify arguments in open propositions. Therefore, the different non-canonical strategies that encode subject focus would serve the purpose of \say{detopicalization} of the subject, i.e., they would enhance comprehension of the subject constituent as new information to the hearer, rather than as the entity the sentence is about \parencite{reinhart1981pragmatics}.

Following from these considerations, it might be that a subject in focus conveys an interpretation that is more costly to accommodate for the hearer because of its  unexpectedness. This would also explain why object focus rarely shows a specific formal marking. 

Additionally, the default domain of focus is generally assumed  to be the entire VP \parencite{drubig2003toward}, where objects are generally included. Since a direct object is already found in the default focus domain, additional encoding is not required, and speakers opt for more economic ways, such as canonical structures. 

As has been pointed out, the focus asymmetry is also observable in the majority of Romance languages. Elicitation tasks of Canadian French \parencite{skopeteas2010focus} have shown that speakers overwhelmingly chose to use a cleft sentence to answer a question on the subject, while they opted for the SVO canonical structure whenever the focus was on the object. In the experiment, constructions with clefted objects did not occur at all.\footnote{It must be noted, however, that in this study, speakers were forced to provide full sentences as answers. I thank Andreas Dufter for this suggestion.} 
Thus, to focus non-subject arguments like object, which are generally found sentence-finally, the \textit{in situ} strategy is preferred because it is more economic than using a cleft \parencite{destruel2012french}. 

Economy is said to determine the speaker’s choice among existing structural alternatives for the expression of the same propositional content. This reflects the ‘least effort’ principle \parencite{zipf1949human}, which has been shown to account for several properties of language processing \parencite{skopeteas2010focus, SkopeteasFanselow2011}.
In particular, \textit{in situ} focus does not involve any syntactic operation, and hence qualifies as the least complex structure. Since a cleft construction contains additional structural material, it is plausible to assume that clefting involves a higher degree of structural complexity, and is therefore used only when strictly needed (e.g., in the case of a contrastive or corrective interpretation of the focused object). This generalization, however, does not account for the fact that, even for subject focus, we may also find more economic strategies in French or Spanish. If a cleft is an ideal strategy to focus the subject in French, then why do we find also \textit{in situ} focusing for subjects? Conversely, if a postverbal subject is the means to encode subject focus in Spanish, why do we also find clefts? 
These questions are addressed in the next chapter, where I compare various studies on different focusing strategies for subjects in both languages from a language-internal and a comparative perspective.


